\exercisesection{3}

\begin{enumerate}
\def\labelenumi{\arabic{enumi}.}
\tightlist
\item
  \begin{enumerate}
  \def\labelenumii{(\alph{enumii})}
  \tightlist
  \item
  为使 A-模的族 \((\mathbf{E}_{\lambda})\) 的直和是忠实平坦的，只需每个 \(E_{\lambda}\) 都平坦且其中至少一个是忠实平坦的。
  \end{enumerate}
\end{enumerate}

\begin{enumerate}
\def\labelenumi{(\alph{enumi})}
\setcounter{enumi}{1}
\tightlist
\item
由 (a) 推出：若 A 是单环，则每个非空 A-模都是忠实平坦的。这个结果对半单环成立吗？
\end{enumerate}

\begin{enumerate}
\def\labelenumi{\arabic{enumi}.}
\setcounter{enumi}{1}
\item
设 \((p_n)\) 是素数的严格递增序列，A 是乘积环 \(\prod_{n} \mathbf{Z} / p_n \mathbf{Z}\) 。证明诸 \(\mathbf{Z} / p_n \mathbf{Z}\) 的直和 E 是忠实的投射 A-模，但不是忠实平坦的（注意 E 是 A 的理想且 \(\mathrm{E}^2 = \mathrm{E}\) ）。
\item
设 A 是右凝聚环（§2 习题 12）。为使左 A-模的乘积是忠实平坦的，只需其中每个都平坦且至少一个是忠实平坦的。
\end{enumerate}

由此推出：若 A 是凝聚交换环，则形式幂级数环 \(A[[X_{1},\ldots,X_{n}]]\) 是忠实平坦 A-模。

\begin{enumerate}
\def\labelenumi{\arabic{enumi}.}
\setcounter{enumi}{3}
\item
设 A 是交换域 K 上的单代数，B 是 A 的半单但非单的子代数。证明 A 是忠实平坦的（右或左）B-模，但存在不是忠实平坦的右 B-模 E，而 \(E \otimes_{B} A\) 总是忠实平坦的（习题 1）。
\item
设 A 是交换环，M 是平坦 A-模，它含有非平坦模的子模 N（参看 §2 习题 3）。设 B（相应地 C）是 A-模 A ⊕ N（相应地 A ⊕ M），其中乘法由 \((a, x)(a', x') = (aa', ax' + a'x)\) 定义；则 B 不是平坦 A-模，但 B-模 C 是忠实平坦 A-模，因而 B 满足第 5 目命题 8 的条件。
\item
给出整环 A 与以 A 为子环的环 B 的一个例子，使 B 是平坦 A-模，但存在既非投射又非有限生成的 A-模 E，对此 \(B \otimes_{A} E\) 是有限生成的自由 B-模。
\item
若 K 是域，则环 K{[}X{]} 与域 K(X) 都是忠实平坦 K-模，但 K(X) 不是忠实平坦 K{[}X{]}-模。
\item
设 \(p\) 是素数， \(\mathbf{A}\) 是 \(\mathbf{Q}\) 的由分式 \(k / p^n\) （其中 \(k \in \mathbf{Z}\) ， \(n \geqslant 0\) ）组成的子环。证明 \(\mathbf{A}\) 是平坦 \(\mathbf{Z}\) -模，且存在非平坦的 \(\mathbf{Z}\) -模 \(\mathbf{E}\) ，但 \(\mathbf{A} \otimes_{\mathbf{Z}} \mathbf{E}\) 是平坦 \(\mathbf{A}\) -模。
\item
设 A 是交换环，B 是 A-代数， \((\mathbf{C}_{\lambda})_{\lambda \in L}\) 是 A-代数族，且 \(B_{\lambda} = C_{\lambda} \otimes_{A} B\) 是 \(C_{\lambda}\) 与 B 的张量积代数（对所有 \(\lambda \in L\) ）。设 E 是左 B-模。置 \(E_{\lambda} = B_{\lambda} \otimes_{B} E = C_{\lambda} \otimes_{A} E\) ；这是一个 \((\mathbf{B}_{\lambda}, \mathbf{C}_{\lambda})\) -双模。类似地，若 F 是右 B-模，置
\end{enumerate}

\[
\mathbf {F} _ {\lambda} = \mathbf {F} \otimes_ {\mathbf {B}} \mathbf {B} _ {\lambda} = \mathbf {F} \otimes_ {\mathbf {A}} \mathbf {C} _ {\lambda},
\]

它是一个 \((\mathbf{C}_{\lambda},\mathbf{B}_{\lambda})\) -双模。

\begin{enumerate}
\def\labelenumi{(\alph{enumi})}
\item
证明 \((\mathbf{C}_{\lambda},\mathbf{C}_{\lambda})\) -双模 \(\mathbf{F}_{\lambda}\otimes_{\mathbf{B}_{\lambda}}\mathbf{E}_{\lambda}\) 同构于 \((\mathbf{F}\otimes_{\mathbf{B}}\mathbf{E})\otimes_{\mathbf{A}}\mathbf{C}_{\lambda}\) 。
\item
证明：若 E 是平坦（相应地忠实平坦）B-模，则每个 \(E_{\lambda}\) 都是平坦（相应地忠实平坦）\(B_{\lambda}\) -模。若进一步假设 \(\bigoplus C_{\lambda}\) 是忠实平坦 A-模，则逆命题成立。
\item
证明：若 \(L\) 有限，每个 \(E_{\lambda}\) 都是有限生成投射 \(B_{\lambda}\) -模，且 \(\bigoplus_{\lambda \in L} C_{\lambda}\) 是忠实平坦 A-模，则 \(E\) 是有限生成投射 B-模（利用第 6 目的命题 12）。
\end{enumerate}

\begin{enumerate}
\def\labelenumi{\arabic{enumi}.}
\setcounter{enumi}{9}
\tightlist
\item
  \begin{enumerate}
  \def\labelenumii{(\alph{enumii})}
  \tightlist
  \item
  设 \(\rho: A \to B\) 是环同态。证明：对每个左理想 \(a\) ——属于 \(A\) ，且是子集 \(M\) （属于 \(A\) ）的左零化子——有 \(\rho^{-1}(Ba) = a\) 。
  \end{enumerate}
\end{enumerate}

\begin{enumerate}
\def\labelenumi{(\alph{enumi})}
\setcounter{enumi}{1}
\tightlist
\item
由 (a) 导出同态 \(\rho: A \to B\) 的一个例子，使右 A-模 \(B\) 不是平坦的，但 \(\bar{\rho}^{-1}(Ba) = a\) 对每个左理想 \(a\) （属于 \(A\) ）都成立（参看 §2 习题 17、《代数学》第八章 §3 习题 11、§2 习题 6 以及第九章 §2 习题 4）。
\end{enumerate}

